\section{VisualWebArena：面向真实任务的评估基准}

\subsection{从 WebArena 到 VisualWebArena}
WebArena 已经把网页交互任务结构化，但纯文本表示有明显局限：很多视觉语义（布局、图片、图表、按钮位置）无法完整保留。VisualWebArena 则把视觉信息显式纳入，要求 Agent 同时理解页面截图和结构化网页元素。

\begin{importantbox}{VisualWebArena 解决了什么}
\begin{itemize}
    \item 任务不再只依赖 HTML 文本，而是依赖真实视觉上下文；
    \item 评估不再只测 ``能否写出文本答案''，而测 ``能否完成一段交互轨迹''；
    \item 错误来源更接近真实系统，包括误点、误读、漏检和回溯失败。
\end{itemize}
\end{importantbox}

\begin{knowledgebox}{任务类型通常包含三类}
\begin{enumerate}
    \item \textbf{视觉检索型}：根据图片或版式定位目标；
    \item \textbf{跨站流程型}：多标签页、多网站串联完成任务；
    \item \textbf{条件满足型}：满足价格、评分、时间等约束并完成提交。
\end{enumerate}
\end{knowledgebox}

\subsection{为什么 ``人类容易'' 的任务对 Agent 仍然难}
讲者提到很多任务对人类并不复杂，但对 Agent 依然困难，核心原因是 ``决策链条过长且错误可累积''。一个看似小失误（点错按钮、读错页面元素）就会把后续轨迹全部带偏。

\begin{warningbox}{评估指标的陷阱}
如果只看终局成功率，容易忽略中间失败模式；如果只看中间步骤，又可能高估最终可交付能力。更稳妥的方法是 ``终局指标 + 轨迹质量 + 错误类型'' 联合评估。
\end{warningbox}

\begin{table}[H]
\centering
\small
\begin{tabular}{p{2.8cm}p{4.0cm}p{4.2cm}}
\toprule
\textbf{评估维度} & \textbf{单轮 LLM 常见做法} & \textbf{VisualWebArena 风格做法} \\
\midrule
输入形式 & 纯文本 prompt & 页面截图 + 结构化网页状态 \\
输出形式 & 文本答案 & 多步动作序列 \\
正确性判定 & 文本匹配 & 任务终局 + 轨迹一致性 \\
主要失败形态 & 幻觉/遗漏 & 误操作/回溯失败/策略偏航 \\
\bottomrule
\end{tabular}
\caption{从文本问答评测到交互任务评测的范式变化}
\end{table}

\subsection{本章小结}
VisualWebArena 的价值在于把评估目标从 ``会说'' 变为 ``会做''，并把多模态感知、长链路决策和执行稳定性纳入统一框架。

